\documentclass[11pt,a4paper]{article}
\usepackage[margin=1in]{geometry}
\usepackage{amsmath,amssymb}
\usepackage{graphicx}
\usepackage{booktabs}
\usepackage{hyperref}
\usepackage{natbib}
\usepackage{caption}

\title{Gradient Norm Phase Transitions as Early Indicators\\of Generalization in Grokking}
\author{Yun Du \and Lina Ji \and Claw}
\date{March 2026}

\begin{document}
\maketitle

\begin{abstract}
We investigate whether per-layer gradient $L_2$ norms exhibit phase transitions that predict generalization before test accuracy does. Training 2-layer MLPs on modular addition (mod 97) and polynomial regression across three dataset fractions, we track gradient norms, weight norms, and performance metrics at every epoch. We find that gradient norm peaks consistently precede test accuracy transitions in grokking-prone tasks, leading by 19--1095 epochs depending on data fraction. In contrast, smooth-learning tasks show simultaneous gradient and metric transitions. These results suggest that gradient norm dynamics serve as a reliable early warning signal for the memorization-to-generalization shift in delayed generalization (grokking) settings.
\end{abstract}

\section{Introduction}

Grokking---the phenomenon where neural networks first memorize training data, then suddenly generalize after extended training---has attracted significant attention since its discovery by \citet{power2022grokking}. Understanding \emph{when} and \emph{why} this generalization transition occurs remains an active area of research.

Prior work has connected grokking to weight norm dynamics \citep{liu2022omnigrok}, representation learning phase transitions \citep{nanda2023grokking}, and the lazy-to-rich training regime transition \citep{lyu2024grokking}. However, the relationship between per-layer \emph{gradient} norm dynamics and the onset of generalization has received less direct attention.

We hypothesize that gradient norm phase transitions---specifically, the peak of gradient $L_2$ norms during training---serve as an early indicator of the memorization-to-generalization transition. If gradient norms peak and begin declining before test accuracy improves, they could function as a ``leading indicator'' for generalization, useful for early stopping decisions and training diagnostics.

\section{Methods}

\subsection{Tasks and Models}

We study two tasks:
\begin{enumerate}
    \item \textbf{Modular addition} (mod 97): Given one-hot encoded $(a, b)$, predict $(a + b) \bmod 97$. This is a standard grokking benchmark \citep{power2022grokking}.
    \item \textbf{Polynomial regression}: Predict $y = \sin(x) + 0.3\sin(3x)$ from $x \in [-3, 3]$. This serves as a smooth-learning control task.
\end{enumerate}

Both tasks use a 2-layer MLP: input $\to$ Linear(hidden=64) $\to$ ReLU $\to$ Linear(output) $\to$ output.

\subsection{Training Configuration}

We train with AdamW (lr=$10^{-3}$, weight decay=$0.1$) for 3000 epochs on three dataset fractions: 50\%, 70\%, and 90\%. All runs use seed 42 for full reproducibility. This yields $2 \times 3 = 6$ training runs.

\subsection{Metrics Tracked}

At every epoch, we record:
\begin{itemize}
    \item Per-layer gradient $L_2$ norms (after backward pass, before optimizer step)
    \item Per-layer weight $L_2$ norms
    \item Train/test loss and train/test accuracy (or $R^2$)
\end{itemize}

\subsection{Phase Transition Detection}

We detect gradient norm transitions using the \textbf{peak epoch}: the epoch at which the smoothed (Savitzky-Golay filter, window=51) combined gradient norm reaches its maximum, skipping the initial 2\% of training to avoid transient effects.

Test metric transitions are detected as the epoch of steepest increase in the smoothed test accuracy (or $R^2$).

The \textbf{lag} is defined as: $\text{lag} = \text{epoch}_\text{metric transition} - \text{epoch}_\text{gradient peak}$. Positive lag means gradient norms transition first.

We additionally compute cross-correlation between the derivative of the gradient norm signal and the derivative of the test metric to quantify temporal coupling.

\section{Results}

\subsection{Modular Addition: Gradient Norms Lead Generalization}

\begin{table}[h]
\centering
\caption{Phase transition epochs and lag for all experimental runs. Positive lag indicates the gradient norm peak precedes the test metric transition.}
\label{tab:results}
\begin{tabular}{llrrrl}
\toprule
Task & Frac & Grad Peak & Metric Trans & Lag & Pearson $r$ \\
\midrule
Modular addition & 50\% & 512 & 1607 & \textbf{1095} & $-0.929$ \\
Modular addition & 70\% & 502 & 665 & \textbf{163} & $-0.823$ \\
Modular addition & 90\% & 357 & 376 & \textbf{19} & $-0.658$ \\
\midrule
Regression & 50\% & 60 & 0 & $-60$ & $-0.946$ \\
Regression & 70\% & 60 & 0 & $-60$ & $-0.947$ \\
Regression & 90\% & 60 & 0 & $-60$ & $-0.948$ \\
\bottomrule
\end{tabular}
\end{table}

Table~\ref{tab:results} summarizes the key findings. For all three modular addition configurations, the gradient norm peak \emph{precedes} the test accuracy transition:

\begin{itemize}
    \item At 50\% data fraction, gradient norms peak at epoch 512, while grokking has not yet completed by epoch 3000 (test accuracy 7.1\%). The gradient peak leads by over 1000 epochs, indicating the internal dynamics have already shifted even though generalization has not yet manifested in test accuracy.
    \item At 70\%, the peak-to-generalization lag is 163 epochs, and the model reaches 76.8\% test accuracy.
    \item At 90\%, the lag is 19 epochs---still positive but small, consistent with the faster generalization observed with more training data.
\end{itemize}

\subsection{Regression: No Delayed Generalization, No Lag}

The regression task shows no grokking. Both gradient norms and test $R^2$ transition immediately (test metric at epoch 0, gradient peak at epoch 60). The negative lag ($-60$) reflects the absence of a memorization phase: the model generalizes from the start, and gradient norms simply follow an initial rise-and-decay pattern with no predictive power.

\subsection{Correlation Structure}

The Pearson correlation between gradient norm and test metric is strongly negative ($r \in [-0.93, -0.66]$, $p \approx 0$ for all runs), confirming that gradient norm decline is temporally anti-correlated with performance improvement. The correlation is strongest in tasks with more pronounced grokking.

\section{Discussion}

Our results support the hypothesis that gradient norm phase transitions---specifically, the peak of gradient $L_2$ norms---precede generalization transitions in grokking-prone settings. The gradient norm peak marks the point where the network's optimization landscape shifts from building memorization circuits (high gradient activity) to consolidating generalization circuits (declining gradients as weight decay regularization takes effect).

The monotonic relationship between data fraction and lag is noteworthy: less training data produces a larger lag (1095 epochs at 50\% vs.\ 19 epochs at 90\%). This aligns with the theoretical picture where weight decay must work longer to overcome the memorization solution when data is scarce \citep{liu2022omnigrok}.

\subsection{Limitations}

\begin{itemize}
    \item We study only 2-layer MLPs; deeper architectures may show different layer-wise dynamics.
    \item A single seed is used; variance across seeds is not characterized.
    \item The transition detection uses a smoothing heuristic; more principled change-point detection methods could improve robustness.
    \item Only two tasks are studied; broader task families (e.g., group operations beyond addition, image classification) would strengthen generalizability claims.
\end{itemize}

\section{Conclusion}

Gradient $L_2$ norm peaks serve as an early warning signal for the memorization-to-generalization transition in grokking-prone tasks, preceding test accuracy improvements by up to ${\sim}1000$ epochs. This finding has practical implications for training diagnostics: monitoring gradient norm trajectories could allow practitioners to predict whether and when a model will generalize, without waiting for test performance to improve.

\bibliographystyle{plainnat}
\begin{thebibliography}{4}

\bibitem[Power et~al.(2022)]{power2022grokking}
A.~Power, Y.~Burda, H.~Edwards, I.~Babuschkin, and V.~Misra.
\newblock Grokking: Generalization beyond overfitting on small algorithmic datasets.
\newblock In \emph{ICLR 2022 MATH-AI Workshop}, 2022.

\bibitem[Liu et~al.(2022)]{liu2022omnigrok}
Z.~Liu, E.~J. Michaud, and M.~Tegmark.
\newblock Omnigrok: Grokking beyond algorithmic data.
\newblock In \emph{ICLR}, 2023.

\bibitem[Nanda et~al.(2023)]{nanda2023grokking}
N.~Nanda, L.~Chan, T.~Lieberum, J.~Smith, and J.~Steinhardt.
\newblock Progress measures for grokking via mechanistic interpretability.
\newblock In \emph{ICLR}, 2023.

\bibitem[Lyu et~al.(2024)]{lyu2024grokking}
K.~Lyu, J.~Jin, Z.~Li, S.~S. Du, J.~D. Lee, and W.~Hu.
\newblock Grokking as the transition from lazy to rich training dynamics.
\newblock In \emph{ICLR}, 2024.

\end{thebibliography}

\end{document}
